\documentclass[a4,11pt]{aleph-notas}
% Se puede ver la documentación aquí:
% https://github.com/alephsub0/LaTeX_aleph-notas

% -- Paquetes adicionales
\usepackage{enumitem}
\usepackage{url}
\usepackage{array}
\usepackage{booktabs}
\usepackage{longtable}
\usepackage{aleph-comandos}
\hypersetup{
    urlcolor=blue,
    linkcolor=blue,
}


% -- Datos
\institucion{Facultad de Ciencias Exactas, Naturales y Ambientales}
\carrera{Catálogo STEM}
\asignatura{Matemática III}
\tema{Plan tema 08: Función implícita}
\autor{Andrés Merino}
\fecha{Periodo 2026-1}

\logouno[0.16\textwidth]{Logos/logoPUCE_04_ac}
\definecolor{colortext}{HTML}{1957d1}
\definecolor{colordef}{HTML}{1957d1}
\fuente{montserrat}


\begin{document}

\encabezado

%%%%%%%%%%%%%%%%%%%%%%%%%%%%%%%%%%%%%%%%
\section*{Resultado de Aprendizaje}
%%%%%%%%%%%%%%%%%%%%%%%%%%%%%%%%%%%%%%%%

%%%%%%%%%%%%%%%%%%%%%%%%%%%%%%%%%%%%%%%%
\subsection*{RdA de la asignatura:}
\begin{itemize}[leftmargin=*]
    \item \textbf{RdA 1:} Identificar conceptos, teoremas y propiedades de funciones vectoriales, derivadas parciales e integral vectorial.
    \item \textbf{RdA 2:} Calcular derivadas parciales e integrales con asistentes computacionales.
\end{itemize}

%%%%%%%%%%%%%%%%%%%%%%%%%%%%%%%%%%%%%%%%
\subsection*{RdA del tema:}
\begin{itemize}[leftmargin=*]
    \item Identificar cuándo una función está definida de manera implícita por una ecuación.
    \item Calcular las derivadas parciales de una función definida implícitamente aplicando el teorema de la función implícita.
    \item Interpretar geométricamente la función implícita en el contexto de curvas y superficies de nivel.
\end{itemize}

%%%%%%%%%%%%%%%%%%%%%%%%%%%%%%%%%%%%%%%%
\section*{Introducción}
%%%%%%%%%%%%%%%%%%%%%%%%%%%%%%%%%%%%%%%%

\paragraph{Pregunta inicial:}
La ecuación $x^2 + y^2 = 1$ describe un círculo, pero no define $y$ como una función de $x$ en todo su dominio. ¿Hay alguna manera de calcular cómo cambia $y$ al variar $x$ sin necesidad de despejar $y$ explícitamente?

%%%%%%%%%%%%%%%%%%%%%%%%%%%%%%%%%%%%%%%%
\section*{Desarrollo}
%%%%%%%%%%%%%%%%%%%%%%%%%%%%%%%%%%%%%%%%

%%%%%%%%%%%%%%%%%%%%%%%%%%%%%%%%%%%%%%%%
\subsection*{Actividad 1: Función definida implícitamente}
Se introduce la noción de función definida de manera implícita y se contrasta con la definición explícita.

\paragraph{¿Cómo lo haremos?}
\begin{itemize}[leftmargin=*]
    \item \textbf{Motivación:} retomar la pregunta inicial con el círculo unitario; discutir por qué en muchas situaciones reales (curvas de nivel, restricciones físicas) las relaciones entre variables no se pueden despejar fácilmente.
    \item \textbf{Clase magistral:} se define la función implícita $f$ respecto a $\hat{f}$; se enuncia el teorema que proporciona la fórmula de las derivadas parciales de $f$ en términos de las de $\hat{f}$; se trabajan ejemplos en $\R^2$ y $\R^3$. Se usa el resumen \href{https://andres-merino.github.io/MatematicaIII-05-N0051/01Resumen/Resumen08.pdf}{Resumen08.pdf}.
    \item \textbf{Implementación computacional:} verificar las derivadas obtenidas implícitamente comparándolas con las obtenidas despejando, usando un asistente computacional.
    \item \textbf{Resolución de ejercicios:} calcular derivadas parciales de funciones definidas implícitamente, usando \href{https://andres-merino.github.io/MatematicaIII-05-N0051/04Ejercicios/Ejercicios08.pdf}{Ejercicios08.pdf}.
\end{itemize}

\paragraph{Verificación de aprendizaje:}
\begin{itemize}[leftmargin=*]
    \item Para $x^2+y^2=1$, ¿cuánto vale $\dfrac{dy}{dx}$ usando la función implícita? ¿Coincide con el resultado al despejar?
    \item Para $x^2+y^2+z^2=1$, ¿cuánto valen $\parcial{z}{x}$ y $\parcial{z}{y}$?
    \item ¿Qué condición sobre $\hat{f}$ garantiza que se puede aplicar el teorema?
\end{itemize}

%%%%%%%%%%%%%%%%%%%%%%%%%%%%%%%%%%%%%%%%
\subsection*{Actividad 2: Aplicación a planos tangentes y curvas de nivel}
Se aplica la función implícita para obtener rectas y planos tangentes a curvas y superficies de nivel.

\paragraph{¿Cómo lo haremos?}
\begin{itemize}[leftmargin=*]
    \item \textbf{Clase magistral:} se muestra cómo la función implícita permite calcular la recta tangente a una curva $\hat{f}(x,y)=0$ y el plano tangente a una superficie $\hat{f}(x,y,z)=0$; se conecta con la propiedad del gradiente como vector normal al conjunto de nivel. Se usa el resumen \href{https://andres-merino.github.io/MatematicaIII-05-N0051/01Resumen/Resumen08.pdf}{Resumen08.pdf}.
    \item \textbf{Resolución de ejercicios:} determinar rectas y planos tangentes a curvas y superficies dadas implícitamente, usando \href{https://andres-merino.github.io/MatematicaIII-05-N0051/04Ejercicios/Ejercicios08.pdf}{Ejercicios08.pdf}.
\end{itemize}

\paragraph{Verificación de aprendizaje:}
\begin{itemize}[leftmargin=*]
    \item ¿Cuál es la recta tangente a $x^3+y^3=2$ en el punto $(1,1)$?
    \item ¿Cuál es el plano tangente a $x^2+y^2-z=0$ en el punto $(1,1,2)$?
    \item ¿Por qué el gradiente de $\hat{f}$ es normal a la superficie $\hat{f}(x,y,z)=0$?
\end{itemize}

%%%%%%%%%%%%%%%%%%%%%%%%%%%%%%%%%%%%%%%%
\section*{Cierre}
%%%%%%%%%%%%%%%%%%%%%%%%%%%%%%%%%%%%%%%%

\paragraph{Tarea:}
Resolver del libro \href{https://catalogobiblioteca.puce.edu.ec/cgi-bin/koha/opac-detail.pl?biblionumber=83405&query_desc=su%3A%22An%C3%A1lisis%20vectorial%22}{Cálculo Vectorial de Colley}: sección 2.6, ejercicios 17--30 (impares).

\paragraph{Pregunta de investigación:}
\begin{enumerate}[leftmargin=*]
    \item ¿Qué establece el teorema de la función implícita en su versión completa y qué condiciones garantizan la existencia local de la función?
    \item ¿En qué áreas de la ingeniería o la física aparecen funciones definidas implícitamente que no pueden despejarse explícitamente?
\end{enumerate}

\paragraph{Para el próximo tema:}
Ver los videos:
\begin{itemize}[leftmargin=*]
    \item \href{https://youtu.be/GRXgv9GCOCo?si=PEnJrz5_Sjy36cPX}{Clase sobre Divergencia y Rotacional}.
    \item \href{https://youtu.be/cMtxUdzhBLg?si=FTWGbNGlKSHpQS6W}{Clase sobre Divergencia y Rotacional 2}.
    \item \href{https://youtu.be/rB83DpBJQsE?si=xTd73OwzR8vpNF-k}{Divergencia y curvatura}.
\end{itemize}


\end{document}
